\documentclass[10pt,a4paper]{amsart}
\usepackage[margin=19mm]{geometry}
\usepackage{amsmath,amssymb,mathtools}
\usepackage[scr=rsfs,cal=euler]{mathalfa}
\usepackage{xfrac}
\usepackage[unicode,hidelinks,bookmarks=false]{hyperref}
\newcommand{\ie}{i.e.\ }
\newcommand{\cf}{cf.\ }
\newcommand{\resp}{respectively\ }
\newcommand{\Subsection}{\subsection}
\providecommand{\nobreakdash}{\nobreak\hbox{-}}
\setlength{\emergencystretch}{2em}
\multlinegap0pt
\usepackage{amssymb}
\usepackage{bm}
\usepackage{mathtools}
\usepackage[shortlabels]{enumitem}
\usepackage{xcolor}
\usepackage{natbib}
\usepackage{tikz}
\newtheorem{lemma}{Lemma}
\usepackage{cleveref}
\crefname{lemma}{Lemma}{Lemmas}
\newcommand{\eps}{\varepsilon}
\newcommand\ol[1]{\overline{#1}}
\newcommand{\su}{\subseteq}
\newcommand\wh[1]{\widehat{#1}}
\title{Disproof of the odd Hadwiger conjecture}
\author{Marcus K\"uhn, Lisa Sauermann, Raphael Steiner, Yuval Wigderson}
\date{Source: arXiv:2512.20392v1 (CC BY 4.0)}
\begin{document}
\maketitle

\noindent\emph{Source and licence.} Disproof of the odd Hadwiger conjecture. Version arXiv:2512.20392v1 (CC BY 4.0), distributed by arXiv under the Creative Commons Attribution 4.0 International licence, http://creativecommons.org/licenses/by/4.0/, as recorded in the arXiv licence metadata for this exact version and shown on the abstract page https://arxiv.org/abs/2512.20392v1. Full licence notice: \texttt{licenses/arxiv-2512.20392-CC-BY-4.0.txt}.\par\smallskip

\noindent\textit{Editorial scope.} Counted unit: the article's lemma edges (source0.tex lines 350--354). The article's complete original proof of it is delivered with it (source0.tex lines 357--443, one \texttt{proof} environment of 87 lines, which the article places away from the statement). No mathematics is written, repaired or re-derived: the statement and the proof are the article's own lines, in the article's own notation, and no retained line is substituted or re-flowed.  No further source-local context is needed: every cross-reference of every retained line resolves inside the two delivered spans.  Nothing was omitted from any retained span, so the package carries no omission record.  No retained line contains a citation, so the wrapper carries no bibliography and no bibliography entry.  No retained line contains a figure environment, an image-inclusion command or drawing code, so no image is delivered and the package declares no asset dependency.  Editorial formatting only, in addition: an AMS \texttt{amsart} preamble that restates the article's own declarations and macros used by the retained lines, namely the environments \texttt{lemma} on separate counters, together with the standard packages those lines need.  The retained environments print in this document as Lemma 1 in order of appearance, whereas the article prints the same items with the numbering of its own full document.  The complete native source of the article as distributed in the arXiv e-print for this exact version is bundled unchanged as \texttt{sources/191563/source0.tex}.
\begin{lemma}\label{lem:respectful yields many edges}
	Let $X\in \{R,B\}$, fix $0<\varepsilon<1$, and let $n$ be sufficiently large with respect to $\eps$. Let $P=\{(u_1,v_1),\dots,(u_{s},v_{s})\}$ be an $s$-pairing in $[n]$ such that $s\geq \eps n$. Consider an outcome $\wh{\pi_X}$ of the random map $\pi_X:[n] \to [m]$ which is $\varepsilon$-respectful of $P$ and such that the event $\mathcal F_X$ holds for $\pi_X=\wh{\pi_X}$.
    
    Let $\Pi$ be the set of all pairs $\big\{\{\wh{\pi_X}(u_i),\wh{\pi_X}(u_j)\}, \{\wh{\pi_X}(v_i),\wh{\pi_X}(v_j)\}\big\}$ for all $i,j\in [s]$ with $\wh{\pi_X}(u_i)\ne \wh{\pi_X}(u_j)$ and $\wh{\pi_X}(v_i)\ne \wh{\pi_X}(v_j)$ and also $\{\wh{\pi_X}(u_i),\wh{\pi_X}(u_j)\}\ne \{\wh{\pi_X}(v_i),\wh{\pi_X}(v_j)\}$. Then we have $|\Pi|\geq \eps^4n^2/6$.
\end{lemma}

\begin{proof}[Proof of \cref{lem:respectful yields many edges}] Without loss of generality let $X=R$ (the case $X=B$ is analogous). Exclusively for this proof, we adopt the notation $\ol v$ to denote $\wh{\pi_R}(v)$. By the assumption that $\wh{\pi_R}:[n] \to [m]$ is $\varepsilon$-respectful of $P$, there exists a subset $I\su [s]$ of size $|I|\geq \eps s\geq \eps^2n$ 
such that $\ol{u_i} \neq \ol{v_i}$ for all $i \in I$, and furthermore the pairs $\{\ol{u_i},\ol{v_i}\}$ for $ i\in I$ are pairwise distinct. 

Note that for any $i,j\in I$ with $\ol{u_i}\ne\ol{u_j}$ and $\ol{v_i}\ne\ol{v_j}$, we also have
$\{\ol{u_i},\ol{u_j}\}\ne \{\ol{v_i},\ol{v_j}\}$.
Indeed, if $\{\ol{u_i},\ol{u_j}\}= \{\ol{v_i},\ol{v_j}\}$, then due to
$\ol{u_i}\ne \ol{v_i}$ and $\ol{u_j}\ne\ol{v_j}$, we would have
$\ol{u_i}=\ol{v_j}$ and $\ol{u_j}=\ol{v_i}$ and hence
$\{\ol{u_i},\ol{v_i}\}=\{\ol{u_j},\ol{v_j}\}$
for the distinct elements $i,j\in I$
(we must have $i\ne j$ because of $\ol{u_i}\ne\ol{u_j}$).
This would be a contradiction.

Recall that $\Pi$ is the set of all pairs $\big\{\{\ol{u_i},\ol{u_j}\}, \{\ol{v_i},\ol{v_j}\}\big\}$ for all $i,j\in [s]$ with $\ol{u_i}\ne \ol{u_j}$ and $\ol{v_i}\ne \ol{v_j}$ and also $\{\ol{u_i},\ol{u_j}\}\ne \{\ol{v_i},\ol{v_j}\}$. Now, let us just consider the pairs
$\{\{\ol{u_i},\ol{u_j}\}, \{\ol{v_i},\ol{v_j}\}\}$
for all $i,j\in I$ with $\ol{u_i}\ne\ol{u_j}$ and $\ol{v_i}\ne\ol{v_j}$ (recall that we then also have $\{\ol{u_i},\ol{u_j}\}\ne \{\ol{v_i},\ol{v_j}\}$). Together, these pairs form some subset $\Pi_I\su \Pi$,
and it suffices to show $|\Pi_I|\geq \eps^4 n^2/6$.

The key claim is that for any given $x,y,z,w\in [m]$, the pair
$\{\{x,y\}, \{z,w\}\}$
can appear as
$\{\{\ol{u_i},\ol{u_j}\}, \{\ol{v_i},\ol{v_j}\}\}$
for at most two different choices of $\{i,j\}\su I$.
Indeed, if
\[
\big\{\{x,y\}, \{z,w\}\big\}
=
\big\{\{\ol{u_i},\ol{u_j}\}, \{\ol{v_i},\ol{v_j}\}\big\},
\]
then the two pairs $\{\ol{u_i},\ol{v_i}\}$ and $\{\ol{u_j},\ol{v_j}\}$
must be either the two pairs $\{x,z\}$ and $\{y,w\}$ or the two pairs
$\{x,w\}$ and $\{y,z\}$.
More precisely,
$\{\{\ol{u_i},\ol{v_i}\}, \{\ol{u_j},\ol{v_j}\}\}$
must be either $\{\{x,z\}, \{y,w\}\}$ or $\{\{x,w\}, \{y,z\}
\}$.
Since the pairs $\{\ol{u_i},\ol{v_i}\}$ are distinct for all $i\in I$,
in either of these two cases there is at most one possibility for
$\{i,j\}\su I$.
Thus, in total, there are indeed at most two choices for $\{i,j\}\su I$.

Consequently, for any $\{i,j\}\su I$ with
$\ol{u_i}\ne\ol{u_j}$ and $\ol{v_i}\ne\ol{v_j}$, the pair
$\{\{\ol{u_i},\ol{u_j}\}, \{\ol{v_i},\ol{v_j}\}\}$
is an element of $\Pi_I$, and every element of $\Pi_I$ is obtained in this
way for at most two choices of $\{i,j\}\su I$.
Therefore we have
\[
|\Pi_I| \geq \frac{ \big| \big\{\{i,j\}\su I : \ol{u_i}\ne\ol{u_j},\, \ol{v_i}\ne\ol{v_j} \big\} \big|
}{2}.
\]
Among the total number $\binom{|I|}{2}$ of pairs $\{i,j\}\su I$,
the number of pairs with $\ol{u_i}=\ol{u_j}$ is at most
$|I|\cdot (2n/m)$.
Indeed, choosing any $j\in I$, in order to have $\ol{u_i}=\ol{u_j}$,
we need to choose $i\in I$ such that $u_i$ is a vertex in the set
$(\wh{\pi_R})^{-1}(x)$ for $x=\ol{u_j}$.
Due to the event $\mathcal F_R$ we have
$|(\wh{\pi_R})^{-1}(x)|\leq 2n/m$, so there are at most $2n/m$ choices for $u_i$
and hence for $i\in I$
(recall that the vertices $u_1,\dots,u_s,v_1,\dots,v_s$ are all distinct).
Similarly, the number of pairs $\{i,j\}\su I$ with
$\ol{v_i}=\ol{v_j}$ is at most $|I|\cdot (2n/m)$.
Thus,
\begin{align*}
\big|
\big\{\{i,j\}\su I :
\ol{u_i}\ne\ol{u_j},\,
\ol{v_i}\ne\ol{v_j}
\big\}
\big|
&\geq
\binom{|I|}{2}
- |I|\cdot \frac{2n}{m}
- |I|\cdot \frac{2n}{m} \\
&=
|I|\left(\frac{|I|-1}{2}-\frac{4n}{m}\right) \\
&\geq
\eps^2 n
\left(\frac{\eps^2 n-1}{2}-\frac{4n}{m}\right)
\geq
\frac{\eps^4 n^2}{3},
\end{align*}
using that $n$ and hence $m$ is sufficiently large with respect to $\eps$.
This implies
$|\Pi|\geq |\Pi_I|\geq \eps^4 n^2/6$, as desired.
\end{proof}

\end{document}
